\documentclass[11pt]{article}
\usepackage[a4paper,margin=25mm]{geometry}
\usepackage[T1]{fontenc}
\usepackage{lmodern,amsmath,amssymb,amsthm,microtype,hyperref}
\title{Addendum: GHS curvature and Riemann theta heat flow}
\author{Tony Newton}
\date{}
\begin{document}\maketitle
The classical Riemann theta potential $V(u)=-\log\Phi(u)$
was studied by Newman, who proved in 1991 that $V'$
is convex on the positive half-line \cite{GHS}.
Thus the qualitative monotone-curvature assertion in
the present preprint is previously known. Our rational
interval certificate should be regarded as an alternative
quantitative verification, not a first proof of the
underlying GHS-type property.

We use the de Bruijn--Newman deformation in the
normalization employed by Rodgers and Tao \cite{RT}:
\[
H_t(z)=\int_0^\infty e^{tu^2}\Phi(u)\cos(zu)\,du .
\]
Newman and Wu review this deformation and its
statistical-mechanical background \cite{NW}. For
$V_t=-\log(e^{tu^2}\Phi)$ we have exactly
\[
V_t''=V_0''-2t,\qquad
V_t^{(k)}=V_0^{(k)}\quad(k\ge3).
\]
Thus every derivative of the potential of order
at least three is independent of the heat time.
Nevertheless known de Bruijn--Newman theory gives
$0\le\Lambda\le1/2$: $H_{1/2}$ has only
real zeros, whereas $H_{-1}$ has nonreal
zeros \cite{RT,NW}. The entire third-and-higher
potential jet cannot determine the real-zero property.

Differentiating the heat Fourier integral gives
$\partial_tH_t=-\partial_z^2H_t$. With
\[
Q_t(x,y)=H_t(x+iy)H_t(x-iy)
=\sum_{n\ge0}L_{n,t}(x)y^{2n}
\]
one obtains
\[
\partial_tQ_t=\tfrac12(Q_{yy}-Q_{xx}),\quad
\partial_tL_{n,t}
=(n+1)(2n+1)L_{n+1,t}
-\tfrac12\partial_x^2L_{n,t}.
\]
The negative adjacent-order source on backward
heat evolution defeats a direct scalar positivity
comparison. An explicit Gaussian--hyperbolic
double-zero example in the cumulative manuscript
shows why one proposed universal comparison
bound is too strong.

The same manuscript provides a theta-specific,
heat-uniform \emph{absolute} Fourier truncation
estimate for $0\le t\le1/2$:
\[
|F_t^{(k)}(z)-F_{t,M}^{(k)}(z)|
\le \frac{3k!}{4^k}q e^{-q+1/16},
\quad q=\pi(M+1)^2,\quad|\Im z|\le3 .
\]
For $k=0$ the bound also holds on $|\Im z|\le7$.
Together with curvature/moment comparisons it
supports a bounded noncentral interval of
Laguerre positivity uniformly in heat time, but
it does not exclude off-critical zeros globally.

This addendum clarifies earlier references.
The qualitative GHS curvature result belongs
to Newman \cite{GHS}; the zero threshold and
its lower bound belong to the earlier
de Bruijn--Newman work and Rodgers--Tao
\cite{RT}; the relevant broader context is
reviewed by Newman--Wu \cite{NW}.
The Riemann hypothesis remains open.

\begin{thebibliography}{3}
\bibitem{RT} B. Rodgers and T. Tao,
The de Bruijn--Newman constant is non-negative,
\emph{Forum Math. Pi} \textbf{8} (2020), e6.
\bibitem{NW} C. M. Newman and W. Wu,
Constants of de Bruijn--Newman type in analytic
number theory and statistical physics,
\emph{Bull. Amer. Math. Soc.} \textbf{57} (2020), 595--614.
\bibitem{GHS} C. M. Newman,
The GHS inequality and the Riemann hypothesis,
\emph{Constr. Approx.} \textbf{7} (1991), 389--399.
\end{thebibliography}
\end{document}